\documentclass[../../../../SoftwareRequirementsSpecification.tex]{subfiles}
\begin{document}
% Richiede le macro definite in src/macros/useCaseMacros.tex
% (incluse nel preambolo di SoftwareRequirementsSpecification.tex)

\ucstart
\begin{xltabular}{\textwidth}{|l|X|}
  \hline
  \ucid{A-03} \\ \hline
  \ucname{Visualizzazione Schede Assegnate} \\ \hline
  \ucactor{Atleta} \\ \hline
  \ucdescription{L'atleta consulta le schede di allenamento che il
    proprio PT gli ha assegnato, fino al dettaglio degli esercizi
    previsti in ciascuna settimana.} \\ \hline
  \ucprecond{L'atleta deve essere autenticato.} \\ \hline
  \ucpostcond{L'atleta visualizza le sessioni di allenamento previste per la
    settimana scelta.} \\ \hline
  \ucmainflow{
    \item L'atleta apre la pagina delle proprie schede.
    \item Il sistema mostra l'elenco delle schede di allenamento
      assegnate all'atleta, con nome, periodo di validità e stato.
      L'elenco comprende le sole assegnazioni confermate (vedi
      Nota 3).
    \item L'atleta preme su una scheda dell'elenco.
    \item Il sistema mostra il dettaglio della scheda assegnata:
      nome, periodo di validità ed elenco delle settimane di
      allenamento. Il sistema evidenzia la settimana corrente.
    \item L'atleta preme su una settimana, scelta tra tutte quelle
      della scheda, comprese quelle già trascorse (vedi Nota 2).
    \item Il sistema mostra le sessioni di allenamento della
      settimana con gli esercizi previsti e i valori prescritti a
      quell'atleta (vedi Nota 1).
  } \\ \hline
  \ucaltflow{
    \textbf{2a. Nessuna scheda assegnata} \newline
    - Il sistema mostra un messaggio ("Nessuna scheda di allenamento
    assegnata"). \newline
    - Il flusso termina.
  } \\ \hline
  \ucnote{
    \textbf{Nota 1:} i valori mostrati al passo 6 sono quelli
    memorizzati sull'assegnazione: gli esercizi e le loro serie/ripetizioni
      vengono dalle
    sessioni di allenamento configurate nel blocco mensile (PT-04), i
      carichi dai valori che il PT ha
    inserito per quell'atleta al momento dell'assegnazione (PT-05),
    con le personalizzazioni eventualmente introdotte in
    seguito (PT-09).\newline
    \textbf{Nota 2:} tutto ciò che questo caso d'uso mostra è in sola
    lettura, comprese le settimane trascorse. A differenza del PT
    (caso d'uso PT-09), l'atleta non modifica mai ciò che gli è
    prescritto.\newline
    \textbf{Nota 3:} le assegnazioni che il PT ha lasciato in bozza
    non compaiono in questo elenco, perché i loro carichi sono
    incompleti e l'atleta non deve allenarsi su una scheda che il PT
    non ha ancora confermato (vedi Nota 5 del caso d'uso PT-05).
  } \\ \hline
\end{xltabular}
\end{document}
